\documentclass{article}
\usepackage[fleqn]{amsmath}
\usepackage{bm}
\usepackage{amssymb}
\usepackage{graphicx}
\usepackage[title]{appendix}
\usepackage{gensymb}
\newcommand{\Rho}{\mathrm{P}}
\title{Smooth layering algorithm}
\author{Matthieu GENDRIN, Stéphane PATEUX, Théo LADUNE \\
	Orange Innovation
	}
\setlength{\parindent}{0cm}

\date{\today}
% Hint: \title{what ever}, \author{who care} and \date{when ever} could stand 
% before or after the \begin{document} command 
% BUT the \maketitle command MUST come AFTER the \begin{document} command! 
\begin{document}

\maketitle

\section{Model and terminology}

Consider the color is given by:
\begin{align}
    \label{model}
    \text{pixel \textbf{C}olor: } & C = \sum_{k} T_k C_k E_k \\
    \text{cumulated \textbf{T}ransmittance: } & T_k = \prod_{p<k} t_p \nonumber \\
    \text{layer \textbf{t}ransmittance: } & t_k = \exp(-A_k (b_{k+1}-b_k)) \nonumber \\
    \text{layer \textbf{A}bsorptivity: } & A_k = \sum_{i \in \mathbb{G}_k} a_i \nonumber \\
    \text{gaussian \textbf{a}bsorptivity: } & a_i = d_i / (b_{f_i}-b_{n_i}) \nonumber \\
    \text{layer \textbf{C}olor: } & C_k = \frac{1}{A_k} \sum_{i \in \mathbb{G}_k} a_i c_i \nonumber \\
    \text{layer \textbf{E}missivity: } & E_k = 1 - t_k \nonumber
\end{align}

\paragraph{Discussion:}
$d_i, c_i, b_{n_i}, b_{n_i}$ are considered here as the input parameters. In the back propagation, one will need to 
backpropagate through $d_i, b_{n_i}, b_{f_i}$ which are functions of the actual 3d parameters:
\begin{align}
    d_i & = \rho_i * \exp(-\frac{uv^T\Sigma_i^{-1}uv}{2}) \nonumber \\
    \Sigma_i & = f(xyz_i, scale_i, rot_i) \text{ as in 3DGS} \nonumber \\
    b_{n_i}, b_{f_i} & = h(xyz_i, scale_i, rot_i) \nonumber
\end{align}
$\rho_i, xyz_i, scale_i, rot_i, c_i$ are initial parameters attached to each gaussian kernel. \\
$n_i$ (as \textbf{n}ear) is the idx of the entry bound of gaussian $i$ in the ordered list of all bounds. \\
$f_i$ (as \textbf{f}ar) is the idx of the exit bound of gaussian $i$ in the ordered list of all bounds.

\section{Derivation / step I}


Derivatives of \ref{model} gives:
\begin{align}
    \frac{\partial C}{\partial a_i} & = b_{n_i} R_{n_i} - b_{f_i} R_{f_i} - Y_{n_i} + Y_{f_i} + c_i (Z_{n_i} - Z_{f_i}) + X_{n_i} - X_{f_i} \\
    \frac{\partial C}{\partial c_i} & = a_i (Z_{n_i} - Z_{f_i}) \\
    \frac{\partial C}{\partial b_{n_i}} & = A_{n} R_n + T_{n-1} C_{n-1} A_{n-1} t_{n-1} - T_n C_n A_n t_n + \frac{d_i}{(b_{f_i} - b_{n_i})^2} \frac{\partial C}{\partial a_i} \\
    \frac{\partial C}{\partial b_{f_i}} & = A_{n} R_n + T_{n-1} C_{n-1} A_{n-1} t_{n-1} - T_n C_n A_n t_n - \frac{d_i}{(b_{f_i} - b_{n_i})^2} \frac{\partial C}{\partial a_i} \\
    \text{with:} \nonumber \\
    R_k & = \sum_{p \geq k}^{L+1} T_p C_p E_p, (C_{L+1}=bg, E_{L+1} = 1) \nonumber \\
    Y_k & = \sum_{p \geq k}^{L} T_p C_p E_p b_p \nonumber \\
    Z_k & = \sum_{p \geq k}^{L} \frac{T_p E_p}{A_p} \nonumber \\
    X_k & = \sum_{p \geq k}^{L} T_p C_p t_p (b_{p+1}-b_p) - \frac{T_p C_p E_p}{A_p}\nonumber
\end{align}
Where $L$ is the index of the last layer ($T_L \geq T_{min}$, $T_{L+1} < T_{min}$). Note that with $T_{min}$ very small, 
the sum is equivalent with the sum over all the layers.

\paragraph{Discussion:}
Pose of the gaussian $i$ only backpropagates through $\frac{\partial C}{\partial b_{n_i}}$ and $\frac{\partial C}{\partial b_{f_i}}$. \\
$\frac{\partial C}{\partial \rho_i}$ and $\frac{\partial C}{\partial c_i}$ can be split into a sum of partial contributions,
the first from $n_i$ to $n$ and the second from $n$ and $f_i$, for any $n \in ]n_i,f_i[$.

\section{Algo proposal for back propagation}

Given the hypothesis that we have a list of contributing gaussian for the current layer. (WARNING: this list is not bounded!). 
We can maintain a number of limited-sized buffers $R, U, V, Y, Z$ of size $N$. Those layer can be updated with the last
$N$ values at each layer, given the hypothesis that we have the final value (processed in forward pass).
At each step (layer $k$):
\begin{enumerate}
    \item the buffers are updated (which involves processing the contribution of all contributing gaussians)
    \item the exiting gaussian backpropagates (if the bound $k$ is an exit) and is removed from the contributing gaussians
    \item the entering gaussian is added to contributing gaussians (if the bound $k$ is an entry)
    \item for any contributing gaussian, if $(k - n_i) \equiv 0 \mod (N - 1)$, then the gradient of C for the parameters of this gaussian are updated with partial contribution for $[k - (N - 1), k]$.
\end{enumerate}


\begin{appendices}
    
    \section{Derivation}
    
    \subsection{Derivation / step I}
    
    As a first step, we'll not consider $a_i$ as an input parameter.

    Intermediate derivations:
    \begin{align}
        \frac{\partial T_k}{\partial t_p} & = \delta_{p<k} \frac{T_k}{t_p} \nonumber \\
        \frac{\partial t_k}{\partial a_i} & = - \delta_{n_i \leq k < f_i} t_k (b_{k+1}-b_k) \nonumber \\
        \frac{\partial t_k}{\partial b_p} & = A_k t_k (\delta_{p,k} - \delta_{p,k+1}) \nonumber \\
        \frac{\partial C_k}{\partial a_i} & = \delta_{n_i \leq k < f_i} ( \frac{-1}{A_k^2} \sum_{j \in \mathbb{G}_k} a_j c_j + \frac{1}{A_k} c_i) = \delta_{n_i \leq k < f_i} \frac{c_i-C_k}{A_k} \nonumber \\
        \frac{\partial C_k}{\partial c_i} & = \delta_{n_i \leq k < f_i} \frac{a_i}{A_k} \nonumber \\
        \frac{\partial E_k}{\partial t_p} & = - \delta_{p,k} \nonumber
    \end{align}

    Deriving color from (\ref{model}) by $a_i$ gives:
    \begin{align}
        \frac{\partial C}{\partial a_i} & = \sum_{k=0}^{L} \frac{\partial T_k}{\partial a_i} C_k E_k \label{dcolor_dai_A} \\
                                        & + \sum_{k=0}^{L} T_k \frac{\partial C_k}{\partial a_i} E_k \label{dcolor_dai_B} \\
                                        & + \sum_{k=0}^{L} T_k C_k \frac{\partial E_k}{\partial a_i} \label{dcolor_dai_C}
    \end{align}
    Part (\ref{dcolor_dai_A}):
    \begin{align}
        \frac{\partial T_k}{\partial a_i} & = 0, \forall k \in [0,n_i] \nonumber \\
                                          & = -T_k \sum_{n_i \leq p < k} (b_{p+1} - b_p) = -T_k (b_k - b_{n_i}), \forall k \in [n_i, f_i[ \nonumber \\
                                          & = -T_k \sum_{n_i \leq p < f_i}(b_{p+1} - b_p) = -T_k (b_{f_i} - b_{n_i}), \forall k \geq f_i \nonumber \\
        \text{thus:} \nonumber \\
        \sum_{k=0}^{L} \frac{\partial T_k}{\partial a_i} C_k E_k & = \sum_{k=n_i}^{min(f_i-1,L)} -T_k (b_k - b_{n_i}) C_k E_k + \sum_{k=f_i}^{L} -T_k (b_{f_i} - b_{n_i}) C_k E_k \nonumber \\
                                                                 & = b_{n_i} \sum_{k=n_i}^{L} T_k C_k E_k - b_{f_i} \sum_{k = f_i}^{L} T_k C_k E_k - \sum_{k=n_i}^{min(f_i-1,L)} T_k C_k E_k b_k \nonumber \\
        \sum_{k=0}^{L} \frac{\partial T_k}{\partial a_i} C_k E_k & = b_{n_i} \sum_{k=n_i}^{L} T_k C_k E_k - b_{f_i} \sum_{k = f_i}^{L} T_k C_k E_k \nonumber \\
                                                                 & - \sum_{k=n_i}^{L} T_k C_k E_k b_k + \sum_{k=f_i}^{L} T_k C_k E_k b_k \nonumber
    \end{align}
    Part (\ref{dcolor_dai_B}):
    \begin{align}
        \sum_{k=0}^{L} T_k \frac{\partial C_k}{\partial a_i} E_k & = \sum_{k=}^{L} T_k \delta_{n_i \leq k < f_i} \frac{c_i-C_k}{A_k} E_k \nonumber \\
                                                                 & = \sum_{k=n_i}{min(f_i-1,L)} T_k \frac{c_i-C_k}{A_k} E_k \nonumber \\
                                                                 & = c_i \sum_{k=n_i}^{min(f_i-1,L)} \frac{T_k E_k}{A_k} - \sum_{k=n_i}^{min(f_i-1,L)} \frac{T_k C_k E_k}{A_k} \nonumber \\
                                                                 & = c_i (\sum_{k=n_i}^{L} \frac{T_k E_k}{A_k} - \sum_{k=f_i}^{L} \frac{T_k E_k}{A_k}) \nonumber \\
                                                                 & - \sum_{k=n_i}^{L} \frac{T_k C_k E_k}{A_k} + \sum_{k=f_i}^{L} \frac{T_k C_k E_k}{A_k} \nonumber
    \end{align}
    Part (\ref{dcolor_dai_C}):
    \begin{align}
        \sum_{k=0}^{L} T_k C_k \frac{\partial E_k}{\partial a_i} & = \sum_{k=0}^{L} T_k C_k (-1) (- \delta_{k=n_i}^{min(f_i-),L} t_k (b_{k+1}-b_k)) \nonumber \\
                                                                 & = \sum_{k=n_i}^{min(f_i-),L} T_k C_k t_k (b_{k+1}-b_k) \nonumber \\
                                                                 & = \sum_{k=n_i}^{L} T_k C_k t_k (b_{k+1}-b_k) - \sum_{k=f_i}^{L} T_k C_k t_k (b_{k+1}-b_k) \nonumber
    \end{align}
    Gatering the three parts:
    \begin{align}
        \frac{\partial C}{\partial a_i} & = b_{n_i} \sum_{k=n_i}^{L} T_k C_k E_k - b_{f_i} \sum_{k = f_i}^{L} T_k C_k E_k \nonumber \\
                                        & - \sum_{k=n_i}^{L} T_k C_k E_k b_k + \sum_{k=f_i}^{L} T_k C_k E_k b_k \nonumber \\
                                        & + c_i (\sum_{k=n_i}^{L} \frac{T_k E_k}{A_k} - \sum_{k=f_i}^{L} \frac{T_k E_k}{A_k}) \nonumber \\
                                        & - \sum_{k=n_i}^{L} \frac{T_k C_k E_k}{A_k} + \sum_{k=f_i}^{L} \frac{T_k C_k E_k}{A_k} \nonumber \\
                                        & + \sum_{k=n_i}^{L} T_k C_k t_k (b_{k+1}-b_k) - \sum_{k=f_i}^{L} T_k C_k t_k (b_{k+1}-b_k) \nonumber \\
        \frac{\partial C}{\partial a_i} & = b_{n_i} R_{n_i} - b_{f_i} R_{f_i} - Y_{n_i} + Y_{f_i} + c_i (Z_{n_i} - Z_{f_i}) + X_{n_i} - X_{f_i} \label{dcolor_dai} \\
        \text{with:} \nonumber \\
        R_k & = \sum_{p \geq k}^{L+1} T_p C_p E_p, (C_{L+1}=bg, E_{L+1} = 1) \nonumber \\
        Y_k & = \sum_{p \geq k}^{L} T_p C_p E_p b_p \nonumber \\
        Z_k & = \sum_{p \geq k}^{L} \frac{T_p E_p}{A_p} \nonumber \\
        X_k & = \sum_{p \geq k}^{L} T_p C_p t_p (b_{p+1}-b_p) - \frac{T_p C_p E_p}{A_p}\nonumber
        \end{align}

    Deriving color from (\ref{model}) by $c_i$ gives:
    \begin{align}
        \frac{\partial C}{\partial c_i} & = \sum_{k} T_k \frac{\partial C_k}{\partial c_i} E_k \nonumber \\
        & = \sum_{k} T_k \delta_{n_i \leq k < f_i} \frac{a_i}{A_k} E_k \nonumber \\
        & = a_i \sum_{n_i \leq k < f_i} \frac{T_k E_k}{A_k} \nonumber \\
        & = a_i (\sum_{k \geq n_i} \frac{T_k E_k}{A_k} - \sum_{k \geq f_i} \frac{T_k E_k}{A_k}) \nonumber \\
        \frac{\partial C}{\partial c_i} & = a_i (Z_{n_i} - Z_{f_i}) \label{dcolor_dci}
    \end{align}

    Deriving color from (\ref{model}) by $b_n$ gives:
    \begin{align}
        \frac{\partial C}{\partial b_n} & = \sum_{k} \frac{\partial T_k}{\partial b_n} C_k E_k \label{dcolor_dbn_A} \\
                                        & + \sum_{k} T_k C_k \frac{\partial E_k}{\partial b_n} \label{dcolor_dbn_B}
    \end{align}
    Part (\ref{dcolor_dbn_A}):
    \begin{align}
        \sum_{k} \frac{\partial T_k}{\partial b_n} C_k E_k & = \sum_{k} \sum_{p} \delta_{p<k} \frac{T_k}{t_p} A_p t_p (\delta_{n,p} - \delta_{n,p+1}) C_k E_k \nonumber \\
        & = A_{n} \sum_{k > n} T_k C_k E_k \nonumber \\
        & = A_{n} R_n \nonumber
    \end{align}
    Part (\ref{dcolor_dbn_B}):
    \begin{align}
        \sum_{k} T_k C_k \frac{\partial E_k}{\partial b_n} & = \sum_{k} T_k C_k \sum_{p} (- \delta_{p,k}) A_p t_p (\delta_{n,p} - \delta_{n,p+1}) \nonumber \\
        & = - \sum_{k} T_k C_k A_k t_k (\delta_{n,k} - \delta_{n,k+1}) \nonumber \\
        & = T_{n-1} C_{n-1} A_{n-1} t_{n-1} - T_n C_n A_n t_n \nonumber
    \end{align}
    Gatering both:
    \begin{align}
        \frac{\partial C}{\partial b_n} & = A_{n} R_n + T_{n-1} C_{n-1} A_{n-1} t_{n-1} - T_n C_n A_n t_n \label{dcolor_dbn}
    \end{align}

    \subsection{Derivation / step II}

    Now we consider $a_i$ as a function of $b_n$.
    \begin{align}
        \frac{\partial C}{\partial b_{n_i}} & = \frac{\partial C^{(I)}}{\partial b_{n_i}} + \frac{\partial C^{(I)}}{\partial a_i} \frac{\partial a_i}{\partial b_{n_i}}
         = \frac{\partial C^{(I)}}{\partial b_{n_i}} + \frac{\partial C^{(I)}}{\partial a_i} \frac{d_i}{(b_{f_i} - b_{n_i})^2} \\
        \frac{\partial C}{\partial b_{f_i}} & = \frac{\partial C^{(I)}}{\partial b_{f_i}} + \frac{\partial C^{(I)}}{\partial a_i} \frac{\partial a_i}{\partial b_{f_i}}
         = \frac{\partial C^{(I)}}{\partial b_{f_i}} + \frac{\partial C^{(I)}}{\partial a_i} \frac{- d_i}{(b_{f_i} - b_{n_i})^2}
    \end{align}
    Where $\frac{\partial C^{(I)}}{\partial b_n}$ and $\frac{\partial C^{(I)}}{\partial a_i}$ are the ones given in (\ref{dcolor_dbn}).

\end{appendices}

\end{document}
